\title{Parameter-Efficient Orthogonal Finetuning via Factorization}

\begin{abstract}
Large foundation models have become prevalent across numerous domains, yet the expense of training them from the ground up is prohibitive. Consequently, developing methods for efficient adaptation of these models to specific tasks has become crucial. This work investigates Orthogonal Finetuning (OFT) as a systematic approach to downstream model adaptation. Although OFT achieves robust generalization, it still incurs a high number of learnable parameters owing to the inherent dimensionality of orthogonal matrices. To mitigate this limitation, we analyze OFT through the lens of information transmission, leading us to articulate essential criteria for enhanced parameter efficiency. Drawing inspiration from the Cooley-Tukey fast Fourier transform algorithm, which exemplifies streamlined information transfer, we propose a compact orthogonal parameterization leveraging butterfly structures. We incorporate this structure within OFT, resulting in a new, parameter-efficient finetuning approach—Orthogonal Butterfly~(BOFT). BOFT generalizes OFT, subsuming it as a particular instance, and provides a broader framework for orthogonal finetuning. Finally, we present comprehensive experiments demonstrating the adaptation of large vision transformers, large language models, and text-to-image diffusion models to a range of vision and language downstream tasks.
\end{abstract}

\section{Introduction}

Models such as ChatGPT~\cite{brown2020language,chen2021evaluating} and Stable Diffusion~\cite{rombach2022high} have highlighted the exceptional generalization capabilities of large foundation models. However, this advancement is accompanied by a substantial growth in model size—

(\emph{e.g.}, GPT-3 comprises roughly 175B parameters)—making it increasingly impractical to train such models from the beginning. As a result, significant progress in this domain relies on strategies that allow for the adaptation of these models without resorting to complete retraining. Efficient task adaptation, therefore, has become a cornerstone requirement for leveraging foundation models on downstream tasks. There are chiefly three prevailing approaches for such adaptation: \emph{model finetuning}~\cite{hulora2022,zaken2022bitfit,guo2020parameter,raffel2020exploring,qiu2023controlling,zhang2022adaptive,chavan2023one}, which modifies only a subset of model parameters while preserving the original architecture; \emph{adapter tuning}~\cite{rebuffi2017learning,houlsby2019parameter,pfeiffer2020adapterfusion,lin2020exploring,he2021towards}, which introduces new learnable components to the model; and \emph{prompt tuning}~\cite{li2021prefix,lester2021power}, where extra trainable prefix tokens augment the input. Of these strategies, model finetuning stands out due to its simplicity, effectiveness, and, crucially, its avoidance of any additional inference latency.

The core idea underpinning model finetuning is to maintain a degree of similarity between the pretrained and finetuned models, according to specific criteria, to retain the knowledge acquired during pretraining. Common practice in current finetuning strategies involves employing a modest learning rate, which heuristically limits divergence between the initial and updated models. Given that weight matrices in neural networks exhibit inherent structure, a more systematic method seeks to uphold the relational characteristics within these matrices—specifically, preserving the pairwise angular relationships among neurons. This perspective forms the basis of Orthogonal Finetuning~(OFT)~\cite{qiu2023controlling}, a novel finetuning paradigm in which every neuron within a layer undergoes the same orthogonal transformation, thereby guaranteeing that pairwise angles between neurons are maintained throughout finetuning. While OFT has shown notable success in terms of generalizability and convergence when applied to text-to-image diffusion model finetuning~\cite{qiu2023controlling}, its primary drawback lies in the considerable number of trainable parameters required by full orthogonal matrices, especially in high-dimensional settings. To mitigate this, OFT proposes the use of block-diagonal orthogonal matrices, which curtail the parameter count. Nevertheless, this efficiency introduces certain limitations: the imposed sparsity pattern is rigid, and orthogonal transformations are executed separately within each block. As a consequence, although this design achieves parameter reduction, it may introduce adverse inductive biases (such as diminished expressive power and the inability to replicate standard linear transformations), and, crucially, determining an optimal sparsity pattern remains an unresolved challenge.

The central challenge lies in constructing a dense orthogonal matrix without incurring an excessive parameter count. Ordinarily, specifying a $d$-dimensional dense orthogonal matrix would necessitate $\mathcal{O}(d^2)$ parameters, which appears prohibitive. To overcome this, our method innovatively builds a dense orthogonal matrix by composing several factorized sparse matrices. This technique exploits the principle that memory usage can be curtailed by exchanging it for additional computation—specifically, the matrix's representation as a product of sparse factors requires repeated multiplications at every training step. To elucidate this matrix factorization, we recast the problem as one of information transmission. Viewed through this lens, the dense orthogonal matrix’s assembly from sparse factors is akin to transmitting information through a grid-like graph topology. The information transmission framework opens up diverse strategies for sparse matrix factorizations that not only limit parameter count but also preserve enough expressiveness to generate dense outputs.

To promote parameter efficiency within this information transmission setting, we turn to the butterfly network configuration found in the Cooley-Tukey fast Fourier transform algorithm~\cite{cooley1965algorithm}, renowned for its efficient data exchange among nodes~\cite{klasing1994broadcasting}. This architecture gives rise to a distinctive sparse matrix factorization, referred to as butterfly factorization. Utilizing this, we can form a dense $d$-dimensional matrix via the product of $\mathcal{O}(\log d)$ sparse matrices, with each containing merely $\mathcal{O}(d)$ nonzero elements. When each of these sparse components is constrained to be orthogonal, the resulting scheme yields an efficient orthogonal parameterization with only $\mathcal{O}(d\log d)$ parameters—a substantial improvement over traditional approaches. By capitalizing on this efficiency, we introduce Orthogonal Butterfly~(BOFT), a novel fine-tuning method that retains the strengths of OFT as a particular instance, thereby furnishing a general orthogonal fine-tuning framework. Notably, both the block-diagonal and butterfly constructions share an underpinning of sparsity, which in both cases curtails the number of required trainable parameters. Compared to LoRA~\cite{hulora2022}, which employs a low-rank parameterization, our method leverages another principal category of structured matrices—sparse matrices—for parameter efficiency, with both approaches falling within the primary structured matrix frameworks identified by~\cite{candes2011robust}.

In contrast to the block-diagonal design adopted by OFT—which negotiates a balance between the expressive power and regularization—BOFT leverages the butterfly structure to achieve a more gradual interpolation between matrices belonging to the full orthogonal group (enhancing expressivity) and the identity matrix (ensuring regularity). This approach allows us to identify a more compact hypothesis space within the orthogonal group for use in downstream applications. Considering that the butterfly structure is integral to many efficient linear transformations, such as the discrete Fourier and cosine transforms, we contend that our choice of structure naturally embeds an inductive bias that both promotes better generalization and reduces overfitting risks. This perspective is underpinned by the fact that the butterfly structure readily captures the behavior of numerous classic linear transforms, whereas the block-diagonal design in OFT lacks this capability. Our main contributions are as follows:

\begin{itemize}[leftmargin=*,nosep]
    \item We approach the challenge of parameter-efficient orthogonal finetuning by introducing a novel information transmission perspective, which reframes the creation of parameter-efficient dense orthogonal matrices as an information transmission task over a grid-structured graph.
    \item Drawing inspiration from the butterfly arrangements used in the Cooley-Tukey algorithm, we develop the Orthogonal Butterfly, a new method for orthogonal finetuning that emphasizes parameter efficiency within our information transmission paradigm.
    \item We furnish several theoretical explanations that elucidate how BOFT is able to substantially lower the quantity of trainable parameters, all while retaining robust expressivity and training stability. BOFT’s matrix factorization further facilitates an interesting form of weight interpolation (see Figure~\ref{fig:boft_interpolation}).
    \item We pioneer the extension of orthogonal finetuning to a broad range of tasks beyond the domain of controllable text-to-image generation~\cite{qiu2023controlling}, highlighting BOFT’s promise as a universal model finetuning approach. Specifically, we demonstrate the applicability of BOFT across multiple adaptation tasks spanning computer vision and natural language processing, where it delivers substantial improvements over contemporary state-of-the-art techniques in terms of parameter efficiency and generalization performance.
\end{itemize}

\textbf{Parameter-efficient finetuning (PEFT)}. As foundation models (\emph{e.g.}, \cite{brown2020language,radford2021learning,rombach2022high,kirillov2023segment}) continue to scale in size and capability, developing methods to adapt these models to downstream tasks in a way that requires updating only a small fraction of parameters has become an active research area~\cite{treviso2023efficient,ding2023parameter,lialin2023scaling}. Within the broad landscape of PEFT techniques~\cite{houlsby2019parameter,aghajanyan2020intrinsic,hulora2022,edalati2022krona,wang2022adamix,gheini2021cross,zaken2022bitfit,guo2020parameter,sung2021training,ansell2022composable,lester2021power,li2021prefix,vu2022spot,he2021towards,mao2021unipelt,karimi2021compacter,liu2022few,sung2022lst,chen2022parameter,jia2022visual,chen2022adaptformer,zhang2022neural,jie2023fact,lian2022scaling,luo2023towards}, reparameterization-based strategies~\cite{aghajanyan2020intrinsic,hulora2022,edalati2022krona} are most pertinent to our study. The LoRA approach~\cite{hulora2022} enhances a pretrained weight matrix by supplementing it with the product of two trainable low-rank matrices, achieving compelling results, particularly for NLP tasks. In LoRA, the ranks of these added matrices are fixed, whereas later works~\cite{zhang2022adaptive,valipour2022dylora,zhang2023increlora} explore mechanisms to adaptively vary the rank across layers, optimizing parameter usage according to need. The practicality and effectiveness of low-rank reparameterization have made such methods highly influential~\cite{dettmers2023qlora,zi2023delta,chavan2023one}. Motivated by research into hyperspherical energy and its relevance to generalization~\cite{liu2018learning,liu2021orthogonal}, \cite{qiu2023controlling} introduce orthogonal finetuning (OFT)—an alternative method particularly suited to text-to-image diffusion model adaptation. Here, a layer-specific orthogonal transformation is learned, enabling better generalization and more consistent training than LoRA, though typically at the cost of increasing the number of trainable parameters. Consequently, improving the parameter efficiency of OFT remains a meaningful direction. Additionally, the potential of OFT outside the domain of text-to-image diffusion tasks~\cite{qiu2023controlling} has yet to be fully explored. BOFT advances OFT by utilizing butterfly factorization to reduce the number of trainable parameters, thereby expanding the practical applicability of orthogonal finetuning to general adaptation scenarios.

\textbf{Butterfly structures}. The radix-2 Cooley-Tukey algorithm offers an efficient recursive scheme for the $N$-point discrete Fourier transform, decomposing it into two smaller $\frac{N}{2}$-point DFTs. This recursive reduction gives rise to a butterfly structure, formally expressed as the product of several sparse matrices—a composition referred to as a butterfly matrix. Butterfly matrices have long been leveraged to parameterize orthogonal matrices in order to circumvent pivoting during Gaussian elimination and to enhance computational efficiency~\cite{parker1995random}. They have been employed to facilitate stable learning in recurrent neural networks~\cite{jing2017tunable} and for kernel function approximation~\cite{munkhoeva2018quadrature}. Recent advances~\cite{dao2019learning,dao2019kaleidoscope} use butterfly-based parameterizations to discover fast linear transformations. Furthermore, the integration of butterfly matrices has enabled more resource-efficient training of neural networks~\cite{chen2022pixelated,dao2022monarch}. Their utility is further evidenced in contexts such as rapid matrix-vector multiplication~\cite{michielssen1996multilevel,o2010algorithm}, efficient approximations of data-sparse matrices~\cite{li2015butterfly}, and optimization of network transmission protocols~\cite{klasing1994broadcasting,li2003linear,rajasingh2016transmission}. In this work, however, we differentiate ourselves by focusing on using butterfly structures specifically to improve the parameter efficiency of OFT.

\section{An Information Transmission View on Orthogonal Finetuning}

We begin by outlining the foundational concepts underlying OFT. OFT adapts a pretrained weight matrix by expressing the updated weight matrix as the product of a learnable orthogonal matrix and the original, fixed pretrained matrix. This differs from LoRA, which finetunes model weights by adding a low-rank matrix to the pretrained weights; in contrast, OFT employs a multiplicative, rather than additive, orthogonal transformation. To limit the number of additional parameters, LoRA employs a low-rank approach, whereas OFT initially uses a block-diagonal structure for its orthogonal matrices~\cite{qiu2023controlling}, with parameter efficiency increasing as the block size decreases. An illustrative comparison of these approaches is presented in Figure~\ref{fig:comp_lora}. The principal rationale for utilizing orthogonal transformations in weight adaptation is to maintain the relative angles between neuron vectors~\cite{liu2018learning,liu2021orthogonal,qiu2023controlling}, thereby better preserving the semantic information acquired during pretraining. In practice, OFT involves optimizing an orthogonal matrix $\bm{R}\in\mathbb{R}^{d\times d}$ for a given pretrained linear layer $\bm{W}^0\in\mathbb{R}^{d\times n}$, altering the standard forward computation from $\bm{z}=(\bm{W}^0)^\top\bm{x}$ to $\bm{z}=(\bm{R}\bm{W}^0)^\top\bm{x}$, where $\bm{x}\in\mathbb{R}^d$ and $\bm{z}\in\mathbb{R}^n$ denote the input and output vectors, respectively. For a neutral starting point, the orthogonal matrix $\bm{R}$ is initialized as the identity matrix. To guarantee that $\bm{R}$ remains orthogonal during the training process, the Cayley parameterization~\cite{qiu2023controlling,liu2021orthogonal} is adopted: specifically, $\bm{R} = (\bm{I}+\bm{Q})(\bm{I}-\bm{Q})^{-1}$, with $\bm{Q}$ being a skew-symmetric matrix, satisfying $\bm{Q} = -\bm{Q}^\top$. To further reduce parameter count, $\bm{R}$ is constructed in block-diagonal form, parameterized as $\text{diag}(\bm{R}_1, \bm{R}_2, \cdots, \bm{R}_r)$, with each $\bm{R}_i\in\mathcal{R}^{b\times b}$ a smaller orthogonal matrix such that $br = d$. However, this block-diagonal parameterization presupposes that the neuron dimensions (i.e., columns of $\bm{W}^0$) are partitioned into $r$ separate groups, with distinct orthogonal transformations applied to each, which may not correspond to any meaningful decomposition. While block-diagonal orthogonal matrices are effective in practice, splitting neuron dimensions arbitrarily according to their indices lacks theoretical justification, thus making dense orthogonal matrices more appealing. This naturally leads to the following question: \emph{Is it possible to design a dense orthogonal matrix that does not sacrifice parameter-efficiency?}

In order to tackle this problem, we suggest constructing a dense orthogonal matrix by multiplying several sparse orthogonal matrices. To bring clarity and cohesion to the analysis of sparsity patterns within orthogonal matrix factorization, we reinterpret the generation of a dense orthogonal matrix in OFT as an instance of information transmission. Concretely, the process of obtaining a dense matrix $\bm{R}\in\mathbb{R}^{d\times d}$ through the product of $m$ square matrices, written as $\thickmuskip=2mu \medmuskip=2mu \bm{R}=\bm{B}_m\bm{B}_{m-1}\cdots\bm{B}_1$, can be conceptualized as transmitting data across a $d\times (m+1)$ grid of nodes, as depicted in Figure~\ref{fig:info_trans}. The inspiration for this transmission-centric framework arises from the perspective that a dense square matrix of size $d$ encodes complete connections between two groups of $d$ nodes each. Within $\bm{R}$, a zero entry at $R_{ij}$ signifies the impossibility of passing information from node $j$ to node $i$, whereas a nonzero value indicates a viable communication pathway. Thus, the decomposition of the dense matrix $\bm{R}$ into a sequence $\bm{B}_m\bm{B}_{m-1}\cdots\bm{B}_1$ can be viewed as consecutive rounds of information exchange, each governed by the graph induced by the corresponding $\bm{B}_i$. The sequence of transmissions starts with the connections prescribed by $\bm{B}_1$ and culminates with $\bm{B}_n$. To make this concrete, Figure~\ref{fig:info_trans} visualizes an example involving $\bm{R}=\bm{B}_5\bm{B}_4\bm{B}_3\bm{B}_2\bm{B}_1$ along with the respective sparsity patterns and resulting induced graph. Here, the illustrated graph unfolds the process of matrix multiplication, assigning each $\bm{B}_i$ as a matrix dictating the connections from the $i$-th layer of nodes to the $(i+1)$-th. Precisely, the $(j_1,j_2)$ entry in $\bm{B}_i$ determines whether there is a directed edge from node $j_2$ at layer $i$ to node $j_1$ at layer $(i+1)$, with zeros indicating the absence of an edge. For the product $\bm{B}_5\bm{B}_4\bm{B}_3\bm{B}_2\bm{B}_1$ to yield a dense matrix, all nodes in the initial layer must have pathways to every node in the sixth layer. For instance, examining only $\bm{R}=\bm{B}_4\bm{B}_3\bm{B}_2\bm{B}_1$, which focuses on connections between the first and fifth layers, one observes that node $1$ at the start cannot reach node $3$ at the end, which reveals a zero at the $(3,1)$ position in the corresponding sparsity pattern. The analogous relationship between induced graph structure and matrix sparsity remains true for other truncated products such as $\bm{R}=\bm{B}_3\bm{B}_2\bm{B}_1$ and $\bm{R}=\bm{B}_2\bm{B}_1$. More generally, to ensure $\bm{R}\in\mathbb{R}^{d\times d}$ is dense, the sequence of $m$ matrices $\bm{B}_m,\cdots,\bm{B}_1$ must define directed edges on a $d\times (m+1)$ grid, each connecting two adjacent-layer nodes (i.e., between neighboring columns), so that every node at the input layer can transmit information to every node at the output layer $(m+1)$.

Adopting the information transmission perspective offers valuable insight into the sparsity characteristics of the factorization matrices in OFT. In Figure~\ref{fig:blk_diag}, the block-diagonal configuration of $\bm{R}$ in the standard OFT is illustrated. Although this block-diagonal form effectively decreases the number of parameters that must be learned, it is inherently incapable of yielding a fully dense matrix $\bm{R}$. Our objective is thus to generate a dense orthogonal matrix using $m$ sparse orthogonal components, while minimizing the count of effective trainable parameters. Within the information transmission framework, two primary requirements emerge: \emph{(i)} \textbf{dense connectivity}, meaning every initial node should be linked, via some path, to every terminal node; and \emph{(ii)} \textbf{minimum free edges}, where the total edge count remains as small as possible given orthogonality constraints. Notably, orthogonality imposes complex restrictions on connections between consecutive levels. For instance, ensuring that each $\bm{B}_i$ is of full rank—a prerequisite for orthogonality—demands $d$ edges to uniquely pair nodes between levels $i$ and $(i+1)$, guaranteeing at least $d$ connections per adjacent pair of levels (e.g., in Figure~\ref{fig:info_trans}, this is 4). These $d$ obligatory edges should not be included when tallying edge numbers, as they are non-trainable (in a $d \times d$ orthogonal matrix, the $d$ nonzero elements can only be set to $\pm 1$). Owing to their reduced parameterization relative to dense matrices, orthogonal matrices introduce interdependencies among edges; for instance, if a $2\times2$ orthogonal has a zero at the $(1,1)$ location, a corresponding zero must also occur at $(2,2)$—the absence of one edge can necessitate the absence of another. Consequently, any valid edge configuration can be modified by orthogonality constraints, adding or removing edges as needed. Visualizing the sparsity pattern in the resulting orthogonal matrix, this information transmission framework is particularly advantageous in the context of OFT, as our concern lies exclusively with the placement of nonzero, trainable elements in $\bm{R}$, while their specific numerical values are inconsequential.

A straightforward dense linkage between two hierarchical layers involves $\mathcal{O}(d^2)$ edges (that is, the connections described by a single dense orthogonal matrix), which corresponds to $d^2 - d$ individual edges (for instance, with $d=4$, this amounts to 12 edges). Figure~\ref{fig:info_trans} demonstrates a realizable matrix factorization that requires only 10 edges, which is fewer than what is needed for a dense orthogonal matrix. Through this representation, one can analyze the parameter efficiency of OFT from a graph-theoretic viewpoint, making it straightforward to design feasible matrix factorizations. Drawing on the topology of butterfly graphs as employed in the Cooley-Tukey algorithm~\cite{cooley1965algorithm}, this approach permits the efficient full interconnection between $d$ input nodes and $d$ output nodes, utilizing only $\mathcal{O}(d\log d)$ edges. In particular, the configuration illustrated in Figure~\ref{fig:info_trans} requires 10 edges for dense connectivity, whereas the butterfly network achieves the same with only 8 edges. The following section discusses how the butterfly architecture yields improvements in parameter efficiency.

\section{Orthogonal Parameterization by Butterfly Factorization}

Initially introduced in the Cooley-Tukey algorithm for the computation of fast Fourier transforms, the butterfly structure serves to efficiently distribute local modifications in the frequency domain to global modifications in the spatial domain—a property closely aligned with the problem of information dissemination in our context, where each node at the initial layer communicates with all nodes at the terminal layer. Moreover, the butterfly arrangement is widely adopted as an efficient topology in computer networks~\cite{solihin2015fundamentals,li2011linear}. Assuming that $k\geq 2$ and $k$ is a power of $2$, we formally define the \emph{butterfly factor} $\bm{B}^F(k)$ as follows:

\begin{equation}\label{eq:but_factor}
\begin{aligned}
    \bm{B}^F(k)=\begin{bmatrix}
    \text{diag}(\bm{d}_1) & \text{diag}(\bm{d}_2) \\
    \text{diag}(\bm{d}_3) & \text{diag}(\bm{d}_4)
    \end{bmatrix}\in\mathbb{R}^{k\times k},
\end{aligned}
\end{equation}

where each vector $\bm{d}_i\in\mathbb{R}^{\frac{k}{2}}$ for all $i$. For the case $d=2^N$, the $d$-dimensional \emph{butterfly matrix} $\bm{B}(d)\in\mathbb{R}^{d\times d}$ is recursively characterized as

\begin{equation}
\begin{aligned}
    \!\!\bm{B}(d)=\tilde{\bm{B}}(d,d)\!\cdot\!\begin{bmatrix}
    \bm{B}_1(\frac{d}{2})\!\!\!\!\! & \bm{0} \\
    \bm{0}\!\!\!\!\! &\bm{B}_2(\frac{d}{2})
    \end{bmatrix}= \tilde{\bm{B}}(d,d)\tilde{\bm{B}}(d,\frac{d}{2})\cdots\tilde{\bm{B}}(d,2),
\end{aligned}
\end{equation}

Here, $\bm{B}_1(\frac{d}{2})$ and $\bm{B}_2(\frac{d}{2})$ denote butterfly matrices with dimensionality $\frac{d}{2}$. Next, we introduce the \emph{butterfly component}, expressed as $\thickmuskip=2mu \medmuskip=2mu \tilde{\bm{B}}(d,k)=\text{diag}(\bm{B}^F_1(k),\cdots,\bm{B}^F_{d/k}(k))$, a $d\times d$ block-diagonal matrix composed of blocks of size $k$. In this construction, each block $\bm{B}^F_i(k)$ is a butterfly factor as described in Equation~\ref{eq:but_factor}. With this setup, butterfly matrices can serve as a parameterization for orthogonal matrices, provided that each multiplicative component $\tilde{\bm{B}}(d,k)$ in $\bm{B}(d)$ is orthogonal. Consider initially the block-diagonal matrix $\tilde{\bm{B}}(d,2)$, where blocks are of size $2$; each block can be readily made orthogonal via the Cayley transform (equivalently, a two-dimensional rotation), following approaches in \cite{liu2021orthogonal,qiu2023controlling}. The structure of non-zeros in any butterfly component corresponds to a permutation of that found in $\tilde{\bm{B}}(d,2)$, enabling straightforward orthogonal parameterizations for all butterfly components. This mechanism enables a compact orthogonal matrix parameterization relying on repeated $2\times 2$ orthogonal matrices. Building on \cite{chen2022pixelated}, we further generalize butterfly matrices to define a block butterfly component $\tilde{\bm{B}}^b(d,k)$, in which each element in $\bm{d}_i$ becomes a $b\times b$ matrix. To maintain orthogonality in block butterfly components like $\tilde{\bm{B}}^b(d,2)$, every $2b\times 2b$ block is parametrized as orthogonal. Similarly, the nonzero arrangement in $\tilde{\bm{B}}^b(d,k)$ for $k>2$ is obtained via block-wise permutations of the pattern in $\tilde{\bm{B}}^b(d,2)$, making it possible to ensure orthogonality for these blocks in an analogous manner. Collectively, these constructions define the forward computation in BOFT as

\begin{equation}
    \begin{aligned}\nonumber
    \bm{z}=\big(\bm{R}(m,b)\cdot\bm{W}^0\big)^\top\bm{x},~~~~\text{s.t.}~\bigg{\{}\bm{R}(m,b)=\prod_{i=1}^m \tilde{\bm{B}}^b_{(d,i)}~~\&~~\underbrace{\big(\tilde{\bm{B}}^b_{(d,j)}\big)^\top\tilde{\bm{B}}^b_{(d,j)}=\tilde{\bm{B}}^b_{(d,j)}\big(\tilde{\bm{B}}^b_{(d,j)}\big)^\top\!=\bm{I}_d}_{\forall j\in [1, m]}\bigg{\}},
    \end{aligned}
\end{equation}

Here, for notational convenience, $\tilde{\bm{B}}^b(d,2^{m - i+1})$ is written as $\tilde{\bm{B}}^b_{(d,i)}$, and $\bm{I}_d$ represents the $d$-dimensional identity matrix. The orthogonal transformation $\bm{R}(m,b)\in\mathbb{R}^{d\times d}$ is obtained by sequentially multiplying multiple orthogonal butterfly matrices. To streamline notation, BOFT utilizing $\bm{R}(m,\frac{b}{2})$ is denoted as BOFT($m$,$b$), provided $b\geq 2$. In the specific case where $m=1$, BOFT($1$,$b$) simplifies to a block-diagonal OFT~\cite{qiu2023controlling} with blocks of size $b$. Similarly, when $b$ is set equal to $d$ (i.e., BOFT($1$,$d$)), the formulation matches the original OFT~\cite{qiu2023controlling}, representing a general full orthogonal matrix without block constraints. If the block butterfly matrix $\bm{B}^{\frac{b}{2}}(d)$ is chosen for $\bm{R}$, that is, for BOFT($\log \frac{2d}{b}$,$b$), the result is a dense orthogonal transformation. Generally, BOFT($m$,$b$) introduces $\frac{1}{2}(b-1)dm$ effective, trainable parameters for the fine-tuning of a linear operator of dimensions $d\times n$. Adopting a pure butterfly architecture (i.e., $m=\log d$, $b=2$), the parameter count is reduced to $\mathcal{O}(d\log d)$. This contrasts with the original OFT utilizing a dense orthogonal parameterization, which requires $\mathcal{O}(d^2)$ parameters, and with block-diagonal OFT using $r$ blocks, which employs $\mathcal{O}(bd)$. Consequently, while dense orthogonality in the original OFT mandates $b=d$, the BOFT framework enables arbitrary $b$ to obtain dense orthogonal matrices.

\textbf{Identity initialization for BOFT}. Standard finetuning practices start from pretrained weights to ensure minimal divergence between the updated and original models. Analogous to how LoRA initializes its low-rank matrices to zero, BOFT initializes each butterfly matrix to the identity (meaning the associated skew-symmetric matrix is set to all zeros when applying the Cayley transform).

\textbf{Multiplicative Dropout for BOFT}. Dropout as adopted in LoRA~\cite{hulora2022} guards against overfitting by stochastically deactivating weight updates. While the conventional Dropout mechanism~\cite{srivastava2014dropout} directly applies to LoRA, it does not translate to BOFT due to its multiplicative reparameterization. To resolve this, we introduce a multiplicative Dropout tailored for BOFT: since $\bm{R}(m,b)$ is constructed from $m$ butterfly matrices, each of which can be viewed as $2b\times 2b$ block-diagonal orthogonal structures, the dropout process operates by randomly selecting a fraction $p_1$ of butterfly matrices and a fraction $p_2$ of diagonal blocks within each. Selected blocks are then set to the identity, effectively regularizing the multiplicative structure. 

\section{Intriguing Insights and Discussions}

\textbf{Expressivity of BOFT}. The butterfly-based architecture, combined with permutations, is capable of exactly reproducing numerous well-known fast linear transforms~\cite{dao2019learning,dao2019kaleidoscope} (\emph{e.g.}, fast Fourier transform, Hadamard transform). However, the extent to which our orthogonal butterfly matrix can approximate arbitrary orthogonal matrices remains an open question. To investigate this, we conducted experiments where a randomly sampled dense orthogonal matrix~\cite{anderson1987generation} of dimensions $512\times 512$ was approximated. Figure~\ref{fig:para_eff} shows outcomes averaged across 10 independent random initializations. Here, the y-axis indicates the approximation error, while the x-axis measures the effective count of trainable parameters. Curves of identical colors correspond to BOFT instances sharing the same block size, and the leftmost data point represents the error for BOFT($1$,$b$), i.e., the conventional block-diagonal OFT with block size $b$. Empirically, BOFT demonstrates improved parameter efficiency relative to OFT. For instance, BOFT($9$,$2$) achieves superior expressivity compared to BOFT($1$,$16$) while utilizing far fewer parameters. Typically, configurations of BOFT with smaller $b$ values and higher $m$ values offer better parameter efficiency. For example, BOFT($6$,$4$) requires fewer parameters to reach an approximation error similar to that of BOFT($2$,$16$). Overall, the butterfly matrix yields a more structured subset of the orthogonal group than block-diagonal matrices, establishing that BOFT is provably more expressive than OFT given the same block size.

\begin{theorem}[Expressivity of BOFT]\label{thm:exp_boft}
    BOFT possesses greater expressivity compared to OFT with identical block size. To achieve a universal approximation over all orthogonal matrices of size $d$, one may employ the product $\bm{B}_{d-1,1}(d)\bm{B}^\top_{d-1,2}(d)\cdots\bm{B}_{1,1}(d)\bm{B}^\top_{1,2}(d)$, with each $\bm{B}_{i,j}(d)$ denoting a butterfly matrix.
\end{theorem}

Theorem~\ref{thm:exp_boft} reveals a straightforward pathway to generalize BOFT: the final orthogonal matrix can be formulated as $\thickmuskip=2mu \medmuskip=2mu \bm{R}^G(m_1,b_1,m_2,b_2,l)=\bm{R}_{l,1}(m_1,b_1)\bm{R}_{l,2}^T(m_2,b_2)\cdots\bm{R}_{1,1}(m_1,b_1)\bm{R}_{1,2}^T(m_2,b_2)$, with $\bm{R}_{i,j}^T(m,b)$ referring to the orthogonal matrices employed by BOFT. Setting $m_1=m_2=\log d$, $b_1=b_2=2$, and $l=d-1$, the construction $\bm{R}^G(m_1,b_1,m_2,b_2,l)$ can encompass the full orthogonal group. This composition is also referred to as the kaleidoscope hierarchy~\cite{dao2019kaleidoscope}. Nevertheless, increased expressivity does not universally guarantee better outcomes in the context of finetuning—indeed, even approaches with full expressiveness often display subpar performance when fully finetuned. Achieving optimal model generalization thus requires balancing expressivity and regularization. BOFT broadens the space of tunable parameters while incorporating structural inductive biases, facilitating an improved compromise between model flexibility and generalizability.

\textbf{Spectral properties}. Orthogonal finetuning consistently achieves superior spectral properties compared to LoRA, primarily because it maintains the spectral norm of the pretrained weight matrix $\bm{W}^0$ without alteration. To understand this, consider the singular value decomposition $\bm{W}^0=\bm{U}\bm{\Sigma}\bm{V}^\top$, where $\bm{U}$ and $\bm{V}$ are orthogonal matrices and $\bm{\Sigma}$ contains the singular values along its diagonal. Both OFT and BOFT operate by left-multiplying an orthogonal matrix $\bm{R}$, yielding the updated weights $\bm{R}\bm{U}\bm{\Sigma}\bm{V}^\top$. This operation leaves the principal singular value—that is, the spectral norm of $\bm{W}^0$—unchanged. Preserving the spectral norm has been demonstrated to notably enhance training robustness and generalizability~\cite{miyato2018spectral,yoshida2017spectral}. Additional theoretical insights into these properties are provided in Appendix~\ref{app:sn_property}.

\textbf{Orthogonal finetuning as learning bilinear similarity}. BOFT may also be interpreted as the process of learning a bilinear similarity measure, expressed as $\bm{w}^0_i\bm{R}\bm{x}$, where $\bm{w}^0_i$ denotes the $i$-th column (or neuron) of the original weight matrix $\bm{W}^0$. In this formulation, BOFT is seen as adjusting the similarity matrix $\bm{R}$, subject to orthogonality constraints. This formulation is intimately related to concepts in distance metric learning~\cite{xing2002distance} and the theory of bilinear forms~\cite{roman2005advanced}.

\textbf{Inductive bias and generalization in BOFT}. Because the matrix $\bm{R}(m,b)$ in BOFT typically encodes a well-structured subset of the orthogonal group, it implicitly restricts the space of possible solutions and therefore introduces an inductive bias. We posit that this kind of structured bias, instilled through butterfly factorization, is favorable for generalization, as it captures patterns common to established linear transforms—such as the discrete Fourier, sine/cosine, and Hadamard transforms~\cite{dao2019kaleidoscope}. In addition, the sparse matrix factorization mechanism in BOFT may further confer implicit forms of inductive bias~\cite{gunasekar2017implicit,li2018algorithmic,liu2019neural}.

\textbf{Relation to butterfly-based sparse training}. Several prior studies~\cite{dao2019learning,dao2019kaleidoscope,chen2022pixelated,dao2022monarch} have investigated sparse training approaches utilizing butterfly parameterizations. These efforts generally entail the direct application of the butterfly parameterization to weight matrices, with the neural networks being trained ab initio. The method in~\cite{dao2022monarch} explores the adaptation of pretrained weights by first projecting them onto a specialized class of butterfly matrices, after which the resultant components are optimized for task-specific fine-tuning. In contrast, BOFT adopts a fundamentally different fine-tuning mechanism, in which the model parameters are updated via transformations employing weight matrices that are shared across layers.

\input{rebuttal-results}

\section{Concluding Remarks and Limitations}

This study introduces Orthogonal Butterfly, a broadly applicable, parameter-efficient fine-tuning approach for foundation models that leverages the butterfly matrix structure. The primary innovation underlying our method's parameter efficiency lies in expressing a dense orthogonal matrix as a product of several sparse orthogonal matrices. To systematically discover valid matrix decompositions, we introduce a graphical information transmission perspective. Within this context, we demonstrate that butterfly matrices are particularly effective in fulfilling the requirements of sparse orthogonal matrix factorization. Comprehensive empirical evaluations substantiate the advantages of BOFT for fine-tuning across large language models, vision models, and text-to-image generative systems, further confirming the robustness of BOFT as a universal model adaptation technique.

Notwithstanding its practical performance, BOFT has certain limitations. As the final orthogonal transformation in BOFT results from multiplying several orthogonal matrices, the training process incurs moderately increased computational overhead relative to OFT. Reducing this training time penalty represents a promising direction for future research. On the positive side, after fine-tuning is complete, the learned BOFT orthogonal matrices can be fused into the original pretrained model, so there is no extra latency during inference. Additionally, it is still an open question whether the butterfly structure constitutes the optimal information transmission scheme. Our graphical framework allows for cross-pollination with research from computer networking, a field with a rich literature on evaluating and designing efficient network topologies for information flow. We anticipate that identifying more effective network configurations will lead to improved constructions for dense orthogonal matrix compositions.

\end{document}